\documentclass[11pt,a4paper]{article}
\usepackage[utf8]{inputenc}
\usepackage[T1]{fontenc}
\usepackage[french]{babel}
\usepackage{geometry}
\geometry{margin=2.5cm}
\usepackage{graphicx}
\usepackage{xcolor}
\usepackage{amsmath}
\usepackage{listings}
\usepackage{hyperref}
\usepackage{fancyhdr}
\usepackage{titlesec}
\usepackage{enumitem}
\usepackage{tcolorbox}
\usepackage{float}
\usepackage{tikz}
\usepackage{tabularx}
\usepackage{longtable}
\usetikzlibrary{shapes.geometric,arrows.meta,positioning,calc,fit,backgrounds}

% Configuration des listes
\setlist[itemize]{leftmargin=*, itemsep=3pt, topsep=3pt}
\setlist[enumerate]{leftmargin=*, itemsep=3pt, topsep=3pt}

% Configuration du code
\lstset{
    basicstyle=\ttfamily\small,
    breaklines=true,
    frame=single,
    keywordstyle=\color{blue},
    commentstyle=\color{green!60!black},
    stringstyle=\color{orange},
    tabsize=4,
    showstringspaces=false
}

% En-tête et pied de page
\pagestyle{fancy}
\fancyhf{}
\fancyhead[L]{\small Documentation Technique de l'Algorithme de Similarité - ChatBot SUP'PTIC}
\fancyhead[R]{\small \today}
\fancyfoot[C]{\thepage}

% Environnements personnalisés
\newtcolorbox{important}[1][]{
  colback=red!5!white,
  colframe=red!75!black,
  fonttitle=\bfseries,
  title=#1
}

\newtcolorbox{info}[1][]{
  colback=blue!5!white,
  colframe=blue!75!black,
  fonttitle=\bfseries,
  title=#1
}

\newtcolorbox{astuce}[1][]{
  colback=green!5!white,
  colframe=green!75!black,
  fonttitle=\bfseries,
  title=#1
}

\newtcolorbox{exemple}[1][]{
  colback=orange!5!white,
  colframe=orange!75!black,
  fonttitle=\bfseries,
  title=#1
}

\begin{document}

% Page de titre
\begin{titlepage}
    \centering
    \vspace*{2cm}
    {\Huge \textbf{Documentation Technique de l'Algorithme de Similarité}} 
    
\vspace{0.5cm}
    {\Large ChatBot SUP'PTIC} 

\vspace{1.5cm}
    
    \begin{minipage}{0.8\textwidth}
        \centering
        \small Document de référence pour l'équipe projet \\
        détaillant l'architecture logicielle et le fonctionnement \\
        du moteur de recherche basé sur TF-IDF et Similarité Cosinus
    \end{minipage}
    
\end{titlepage}


% Table des matières
\tableofcontents
\clearpage

% Section 1 : Introduction
\section{Introduction}
\label{sec:introduction}

\subsection{Objectif du Document}

Ce document technique a pour objectif d’expliquer en détail le 
\textbf{fonctionnement de l’algorithme de recherche et de similarité} 
utilisé par le ChatBot SUP'PTIC. 
Il constitue une référence pour les membres de l’équipe projet 
souhaitant comprendre la logique mathématique et le pipeline de traitement 
des questions/réponses (prétraitement, vectorisation TF-IDF, calcul de similarité cosinus).

\begin{important}[Périmètre]
Ce document se concentre \textbf{exclusivement} sur l’algorithme de similarité 
et son intégration dans le moteur de recherche du chatbot. 
\end{important}


\subsection{Ce que Vous Allez Apprendre}

À la fin de ce document, vous serez capable de :

\begin{enumerate}
    \item \textbf{Comprendre l'algorithme} : Saisir comment TF-IDF et la similarité cosinus permettent au chatbot de trouver la bonne réponse
    \item \textbf{Apprécier l'importance des données CSV} : Comprendre pourquoi la qualité de nos fichiers CSV détermine directement la performance du système
    \item \textbf{Contribuer efficacement} : Savoir comment améliorer les données pour optimiser les résultats
    
\end{enumerate}

\subsection{Public Cible}

Ce document s'adresse à :
\begin{itemize}
    \item Les développeurs backend et frontend de l'équipe
    \item Les gestionnaires de contenu responsables de la FAQ
    \item Les designers UX/UI qui conçoivent l'interface
    \item Les nouveaux membres rejoignant le projet
    \item Tout membre de l'équipe souhaitant comprendre le système
\end{itemize}

\begin{astuce}[Prérequis]
Aucune connaissance préalable en intelligence artificielle n'est requise. 
Les concepts mathématiques sont expliqués avec des exemples concrets tirés de notre projet.
\end{astuce}

\clearpage

% Section 2 : L'Algorithme Expliqué
\section{L'Algorithme de Similarité : Cœur du Système}
\label{sec:algorithme}

\subsection{Vue d'Ensemble}

Le chatbot SUP'PTIC repose sur un principe simple mais puissant : \textbf{transformer les phrases en nombres pour pouvoir les comparer mathématiquement}.

\begin{info}[Principe Fondamental]
Quand un utilisateur pose une question, le système :
\begin{enumerate}
    \item Transforme la question en un vecteur de nombres (TF-IDF)
    \item Transforme toutes les questions de la FAQ en vecteurs
    \item Calcule la distance entre la question utilisateur et chaque question FAQ (similarité cosinus)
    \item Retourne la réponse associée à la question FAQ la plus proche
\end{enumerate}
\end{info}

\subsection{TF-IDF : Donner un Poids aux Mots}

\subsubsection{Qu'est-ce que TF-IDF ?}

TF-IDF (Term Frequency - Inverse Document Frequency) est une technique statistique qui évalue l'importance d'un mot dans un document par rapport à un corpus de documents.

\paragraph{TF : Fréquence du Terme}

La fréquence du terme (TF) mesure simplement combien de fois un mot apparaît dans une phrase.

\[
TF(mot, phrase) = \frac{\text{nombre d'occurrences du mot dans la phrase}}{\text{nombre total de mots dans la phrase}}
\]

\paragraph{IDF : Fréquence Inverse du Document}

La fréquence inverse du document (IDF) mesure la rareté d'un mot dans l'ensemble du corpus. Un mot rare est plus discriminant qu'un mot courant.

\[
IDF(mot) = \log \frac{N}{n_t}
\]

où $N$ est le nombre total de documents (questions FAQ) et $n_t$ le nombre de documents contenant le mot.

\paragraph{TF-IDF : La Combinaison}

Le score TF-IDF final combine les deux mesures :

\[
TF\text{-}IDF(mot, phrase) = TF(mot, phrase) \times IDF(mot)
\]

\subsubsection{Exemple Concret avec Notre FAQ}

Prenons un exemple tiré de notre base de données réelle sur les clubs et associations :

\begin{exemple}[Calcul TF-IDF - Exemple Réel]
\textbf{Question utilisateur :} \\
\textit{"Quels sont les frais d'adhésion au club informatique ?"}

\vspace{0.3cm}
\textbf{Questions FAQ candidates :}
\begin{enumerate}
    \item \textit{"Quels sont les documents nécessaires pour adhérer au club informatique ?"}
    \item \textit{"Est-ce que les frais d'adhésion sont remboursables ?"}
\end{enumerate}

\vspace{0.3cm}
\textbf{Analyse des mots clés :}

\begin{itemize}
    \item Le mot \textbf{"frais"} apparaît dans la question utilisateur et dans la FAQ 2, mais \textbf{pas} dans la FAQ 1
    \item Le mot \textbf{"adhésion"} apparaît dans les trois phrases
    \item Le mot \textbf{"club"} apparaît dans la question utilisateur et la FAQ 1
    \item Le mot \textbf{"informatique"} apparaît dans la question utilisateur et la FAQ 1
\end{itemize}

\vspace{0.3cm}
\textbf{Calcul simplifié (supposons 100 questions dans notre base) :}

Pour la question utilisateur "Quels sont les frais d'adhésion au club informatique ?" (8 mots) :

\begin{itemize}
    \item Si "frais" apparaît dans 12 questions : $IDF(\text{frais}) = \log(100/12) \approx 2.12$
    \item $TF(\text{frais}) = 1/8 = 0.125$
    \item $TF\text{-}IDF(\text{frais}) = 0.125 \times 2.12 = 0.27$
\end{itemize}

\begin{itemize}
    \item Si "adhésion" apparaît dans 25 questions : $IDF(\text{adhésion}) = \log(100/25) \approx 1.39$
    \item $TF(\text{adhésion}) = 1/8 = 0.125$
    \item $TF\text{-}IDF(\text{adhésion}) = 0.125 \times 1.39 = 0.17$
\end{itemize}

\begin{itemize}
    \item Si "club" apparaît dans 40 questions : $IDF(\text{club}) = \log(100/40) \approx 0.92$
    \item $TF(\text{club}) = 1/8 = 0.125$
    \item $TF\text{-}IDF(\text{club}) = 0.125 \times 0.92 = 0.12$
\end{itemize}

\vspace{0.3cm}
\textbf{Observation :} Le mot "frais" a le score TF-IDF le plus élevé car il est plus discriminant.
\end{exemple}

\clearpage

\subsection{Similarité Cosinus : Mesurer la Proximité}

\subsubsection{Le Concept Géométrique}

Une fois que nous avons transformé les phrases en vecteurs de nombres, nous devons mesurer leur similarité. La similarité cosinus calcule l'angle entre deux vecteurs dans un espace multidimensionnel.

\[
\text{similarité}(\vec{A}, \vec{B}) = \frac{\vec{A} \cdot \vec{B}}{||\vec{A}|| \times ||\vec{B}||}
\]

où :
\begin{itemize}
    \item $\vec{A} \cdot \vec{B}$ est le produit scalaire (somme des produits des composantes)
    \item $||\vec{A}||$ est la norme du vecteur $\vec{A}$ (racine carrée de la somme des carrés)
    \item $||\vec{B}||$ est la norme du vecteur $\vec{B}$
\end{itemize}

Le résultat varie entre 0 (totalement différent) et 1 (identique).

\subsubsection{Exemple de Calcul Complet}

Reprenons notre exemple précédent avec un calcul complet (simplifié à 3 dimensions pour la clarté).

\begin{exemple}[Calcul de Similarité - Pas à Pas]

Supposons que nous ne considérons que 3 mots : [frais, adhésion, club]

\vspace{0.3cm}
\textbf{Question utilisateur :} \\
\textit{"Quels sont les frais d'adhésion au club informatique ?"}

Vecteur : $\vec{A} = [0.27, 0.17, 0.12]$ (scores TF-IDF calculés précédemment)

\vspace{0.3cm}
\textbf{FAQ 1 :} \\
\textit{"Quels sont les documents nécessaires pour adhérer au club informatique ?"}

Vecteur : $\vec{B_1} = [0.00, 0.18, 0.13]$ (pas de "frais", mais "adhésion" et "club")

\vspace{0.3cm}
\textbf{FAQ 2 :} \\
\textit{"Est-ce que les frais d'adhésion sont remboursables ?"}

Vecteur : $\vec{B_2} = [0.28, 0.19, 0.00]$ ("frais" et "adhésion", mais pas "club")

\textbf{Calcul de similarité($\vec{A}$, $\vec{B_1}$) :}

\begin{enumerate}
    \item Produit scalaire : 
    \begin{align*}
    \vec{A} \cdot \vec{B_1} 
    &= 0.0462
    \end{align*}
    
    \item Norme de $\vec{A}$ :
    \begin{align*}
    ||\vec{A}|| 
    &= \sqrt{0.1162} \approx 0.341
    \end{align*}
    
    \item Norme de $\vec{B_1}$ :
    \begin{align*}
    ||\vec{B_1}|| 
    &= \sqrt{0.0493} \approx 0.222
    \end{align*}
    
    \item Similarité :
    \begin{align*}
    \text{similarité}(\vec{A}, \vec{B_1}) &= \frac{0.0462}{0.341 \times 0.222} \\
    &= \frac{0.0462}{0.0757} \\
    &\approx 0.61 \text{ soit } \mathbf{61\%}
    \end{align*}
\end{enumerate}

\textbf{Calcul de similarité($\vec{A}$, $\vec{B_2}$) :}

\begin{enumerate}
    \item Produit scalaire : 
    \begin{align*}
    \vec{A} \cdot \vec{B_2} 
    &= 0.1079
    \end{align*}
    
    \item Norme de $\vec{B_2}$ :
    \begin{align*}
    ||\vec{B_2}|| 
    &= \sqrt{0.1145} \approx 0.338
    \end{align*}
    
    \item Similarité :
    \begin{align*}
    \text{similarité}(\vec{A}, \vec{B_2}) &= \frac{0.1079}{0.341 \times 0.338} \\
    &= \frac{0.1079}{0.1153} \\
    &\approx 0.94 \text{ soit } \mathbf{94\%}
    \end{align*}
\end{enumerate}

\vspace{0.1cm}
\textbf{Résultat :} Le système sélectionne la FAQ 2 car sa similarité (94\%) est beaucoup plus élevée que celle de la FAQ 1 (61\%).

\end{exemple}

\begin{important}[Conclusion]
La présence du mot clé "frais" dans la FAQ 2 fait toute la différence ! C'est pourquoi la qualité et la richesse du vocabulaire dans nos fichiers CSV sont cruciales.
\end{important}

\subsection{Exemple avec les Programmes Académiques}

Prenons maintenant un exemple tiré de notre FAQ sur les programmes académiques :

\begin{exemple}[Application aux Programmes]

\textbf{Question utilisateur :} \\
\textit{"Combien de temps dure la formation d'ingénieur des télécommunications ?"}

\vspace{0.3cm}
\textbf{Questions FAQ candidates :}
\begin{enumerate}
    \item \textit{"Quelle est la durée de formation pour devenir Ingénieur des Télécommunications ?"}
    \item \textit{"Quelle est la durée de formation pour devenir Technicien des Télécommunications ?"}
    \item \textit{"Quel diplôme obtient-on après le second cycle d'ingénierie ?"}
\end{enumerate}

\vspace{0.3cm}
\textbf{Analyse :}

Mots clés communs :
\begin{itemize}
    \item \textbf{"durée"} et \textbf{"formation"} : présents dans la question utilisateur, FAQ 1 et FAQ 2
    \item \textbf{"Ingénieur"} : présent dans la question utilisateur, FAQ 1 et FAQ 3
    \item \textbf{"Télécommunications"} : présent dans toutes les questions
    \item \textbf{"Technicien"} : uniquement dans FAQ 2 (différent de "Ingénieur")
\end{itemize}

\vspace{0.3cm}
\textbf{Pourquoi FAQ 1 aura le meilleur score :}

\begin{enumerate}
    \item Elle contient \textbf{tous} les mots clés de la question utilisateur
    \item La structure de la phrase est très similaire
    \item Le mot "durée" est un terme discriminant qui apparaît explicitement
    \item Même si FAQ 3 contient "Ingénieur", elle ne mentionne pas "durée" ou "formation"
\end{enumerate}

\vspace{0.3cm}
\textbf{Scores estimés (simplifié) :}
\begin{itemize}
    \item FAQ 1 : $\approx$ 95\% (quasi-identique)
    \item FAQ 2 : $\approx$ 70\% (similaire mais "Technicien" $\neq$ "Ingénieur")
    \item FAQ 3 : $\approx$ 45\% (manque les mots clés "durée" et "formation")
\end{itemize}

\end{exemple}

\begin{astuce}[Leçon Importante]
Plus la question dans notre CSV est formulée de manière proche de ce qu'un utilisateur réel demanderait, meilleur sera le score de similarité. C'est pourquoi nous devons :
\begin{itemize}
    \item Utiliser un vocabulaire naturel et varié
    \item Inclure plusieurs variantes de la même question
    \item Penser comme un étudiant qui pose une question
\end{itemize}
\end{astuce}

\clearpage

% Section 3 : Importance des Données CSV
\section{Importance Cruciale des Données CSV}
\label{sec:donnees_csv}

\subsection{Structure Standard des Fichiers CSV}

Nos fichiers CSV respectent rigoureusement la structure suivante :

\begin{lstlisting}[language=bash, caption=Format CSV Standard]
Question, Reponse, Categorie, SousTheme, Source
\end{lstlisting}

\begin{info}[Colonnes du CSV]
\begin{description}
    \item[Question] : La question telle qu'un utilisateur pourrait la poser naturellement
    \item[Reponse] : La réponse complète, précise et vérifiée
    \item[Categorie] : Catégorie thématique principale (ex: "Vie Étudiante", "Programmes académiques")
    \item[SousTheme] : Sous-thème plus précis (ex: "Clubs et associations", "Présentation des filières")
    \item[Source] : Référence ou contributeur de l'information
\end{description}
\end{info}

\subsection{Pourquoi la Qualité des Questions est Critique}

\subsubsection{Rappel du Principe}

L'algorithme TF-IDF calcule la similarité entre :
\begin{itemize}
    \item La question posée par l'utilisateur
    \item Les questions stockées dans notre fichier CSV
\end{itemize}

\begin{important}[Conséquence Directe]
Si nos questions CSV sont mal formulées, l'algorithme ne trouvera \textbf{jamais} la bonne réponse, même si nous l'avons dans notre base de données !
\end{important}

\subsubsection{Exemples de Bonnes et Mauvaises Questions}

\begin{exemple}[Mauvaise Formulation]

\textbf{❌ Question CSV mal formulée :}

\begin{lstlisting}
Question: FAQ_CLUB_001
Reponse: Pour s'inscrire à un club cette année, il faut...
Categorie: Vie Étudiante
\end{lstlisting}

\textbf{Problème :} 
\begin{itemize}
    \item "FAQ\_CLUB\_001" ne contient aucun mot clé significatif
    \item Si un utilisateur demande "Comment s'inscrire à un club ?", le score TF-IDF sera proche de zéro
    \item Aucune correspondance possible avec une vraie question
\end{itemize}
\end{exemple}

\begin{exemple}[Bonne Formulation - Exemple Réel]

\textbf{✅ Question CSV bien formulée (extrait de notre base réelle) :}

\begin{lstlisting}
Question: Comment s'inscrire à un club cette année ?
Reponse: Pour s'inscrire à un club cette année, il faut se rapprocher du président du club pour se renseigner sur la procédure d'intégration du club.
Categorie: Vie Étudiante
SousTheme: Clubs et associations
Source: Contributeur
\end{lstlisting}

\textbf{Avantages :}
\begin{itemize}
    \item Mots clés pertinents : "inscrire", "club", "année"
    \item Formulation naturelle comme un étudiant poserait la question
    \item Si un utilisateur demande "inscription club" ou "comment rejoindre un club", le score sera élevé
\end{itemize}
\end{exemple}

\clearpage

\subsection{Importance des Variantes de Questions}

Pour maximiser la couverture et la précision, nous devons inclure plusieurs variantes pour chaque réponse.

\begin{exemple}[Variantes - Clubs et Associations]

Pour la même réponse sur l'inscription aux clubs, nous pourrions avoir :

\vspace{0.3cm}
\textbf{Variante 1 (actuelle dans notre base) :}
\begin{lstlisting}
Comment s'inscrire à un club cette année ?
\end{lstlisting}

\textbf{Variante 2 (à ajouter) :}
\begin{lstlisting}
Quelle est la procédure pour rejoindre un club ?
\end{lstlisting}

\textbf{Variante 3 (à ajouter) :}
\begin{lstlisting}
Comment faire pour devenir membre d'une association ?
\end{lstlisting}

\textbf{Variante 4 (à ajouter) :}
\begin{lstlisting}
Inscription club étudiant procédure
\end{lstlisting}

\vspace{0.3cm}
\textbf{Avantages :}
\begin{itemize}
    \item Variante 1 : Couvre "inscrire", "club", "année"
    \item Variante 2 : Couvre "procédure", "rejoindre"
    \item Variante 3 : Couvre "devenir", "membre", "association"
    \item Variante 4 : Couvre les recherches avec mots-clés uniquement
\end{itemize}

Chaque variante augmente les chances qu'un utilisateur trouve la bonne réponse !
\end{exemple}

\subsection{Cas Réels de Notre Base de Données}

Analysons quelques exemples tirés de nos fichiers CSV existants :

\subsubsection{Exemple 1 : Clubs - Question Bien Formulée}

\begin{lstlisting}[caption=Extrait réel - Vie Étudiante]
Question: Est-ce que je peux m'inscrire à plusieurs clubs en même temps ?
Reponse: Oui tu peux t'inscrire à plusieurs clubs au même moment.
Categorie: Vie Étudiante
SousTheme: Clubs et associations
Source: Contributeur
\end{lstlisting}

\begin{astuce}[Pourquoi c'est bien]
\begin{itemize}
    \item Question complète et naturelle
    \item Mots clés : "inscrire", "plusieurs", "clubs", "même temps"
    \item Un utilisateur qui demande "peut-on adhérer à 2 clubs" obtiendra un bon score
    \item Réponse courte et directe
\end{itemize}
\end{astuce}

\subsubsection{Exemple 2 : Programmes - Question Précise}

\begin{lstlisting}[caption=Extrait réel - Programmes académiques]
Question: Quelle est la durée de formation pour devenir Ingénieur des Télécommunications ?
Reponse: La formation dure 2 ans.
Categorie: Programmes académiques
SousTheme: Présentation des filières
Source: À vérifier auprès de SUP'PTIC - https://e-supptic.cm/
\end{lstlisting}

\begin{astuce}[Pourquoi c'est bien]
\begin{itemize}
    \item Question très spécifique avec tous les mots clés nécessaires
    \item "durée", "formation", "Ingénieur", "Télécommunications"
    \item Correspond exactement à ce qu'un candidat demanderait
    \item Réponse concise et factuelle
\end{itemize}
\end{astuce}

\subsubsection{Exemple 3 : Clubs - Question à Améliorer}

\begin{lstlisting}[caption=Question qui pourrait être enrichie]
Question: Quel club choisir si je suis passionné de jeux vidéo ?
Reponse: Le club jeu et loisirs.
Categorie: Vie Étudiante
SousTheme: Clubs et associations
Source: Contributeur
\end{lstlisting}

\begin{important}[Suggestions d'Amélioration]
Cette question est bonne, mais nous pourrions ajouter des variantes :

\begin{itemize}
    \item "Y a-t-il un club pour les amateurs de jeux vidéo ?"
    \item "Club gaming SUP'PTIC existe ?"
    \item "Quel club pour les gamers à l'école ?"
\end{itemize}

Cela augmenterait la couverture pour différentes formulations.
\end{important}

\clearpage

\subsection{Impact du Vocabulaire Riche}

\subsubsection{Synonymes et Variantes Linguistiques}

Les utilisateurs utilisent des mots différents pour exprimer la même idée. Notre CSV doit refléter cette diversité.

\begin{exemple}[Richesse du Vocabulaire]

Pour poser une question sur les frais, un utilisateur peut dire :

\begin{itemize}
    \item "Quels sont les \textbf{frais} d'inscription ?"
    \item "Combien \textbf{coûte} l'inscription ?"
    \item "Quel est le \textbf{montant} à payer ?"
    \item "Quel est le \textbf{prix} de l'inscription ?"
    \item "Quel est le \textbf{tarif} pour s'inscrire ?"
\end{itemize}

\vspace{0.3cm}
\textbf{Mots clés à couvrir :} frais, coûte, coût, montant, prix, tarif, payer, paiement

\vspace{0.3cm}
\textbf{Solution :} Créer plusieurs variantes dans notre CSV utilisant ces différents termes.
\end{exemple}

\subsubsection{Exemple Réel : Diversité du Vocabulaire}

Voici comment notre base actuelle gère différentes formulations :

\begin{lstlisting}[caption=Variantes dans notre base]
# Variante 1
Question: Comment créer un nouveau club à SUP'PTIC ?
Reponse: Pour créer un nouveau club à SUP'PTIC, il faut écrire une demande...

# Variante 2 (à ajouter)
Question: Quelle est la procédure pour lancer une nouvelle association ?
Reponse: [même réponse]

# Variante 3 (à ajouter)
Question: Comment faire pour fonder un club étudiant ?
Reponse: [même réponse]
\end{lstlisting}

\textbf{Vocabulaire couvert :}
\begin{itemize}
    \item Variante 1 : "créer", "nouveau", "club"
    \item Variante 2 : "procédure", "lancer", "association"
    \item Variante 3 : "faire", "fonder", "étudiant"
\end{itemize}

\subsection{Structure CSV et Catégorisation}

\subsubsection{Importance des Métadonnées}

Les colonnes Categorie, SousTheme et Source ne sont pas utilisées directement par l'algorithme de similarité, mais elles sont cruciales pour :

\begin{enumerate}
    \item \textbf{Organisation} : Faciliter la gestion et la maintenance des données
    \item \textbf{Filtrage} : Permettre des recherches par catégorie si nécessaire
    \item \textbf{Validation} : Tracer la source de l'information pour vérification
    \item \textbf{Évolution} : Identifier les zones à enrichir
\end{enumerate}

\subsubsection{Catégories Actuelles}

D'après nos données existantes, nous avons identifié ces catégories principales :

\begin{itemize} \item \textbf{Inscriptions \& Admissions} : procédures d'inscription, conditions d'admission \item \textbf{Examens \& Évaluations} : organisation des examens, modalités d'évaluation \item \textbf{Vie Étudiante} : clubs, associations, activités étudiantes \item \textbf{Services Administratifs} : démarches administratives, services aux étudiants \item \textbf{Programmes Académiques} : présentation des filières, contenus pédagogiques \item \textbf{Infrastructures \& Ressources} : salles, laboratoires, bibliothèque, ressources numériques \item \textbf{Carrières \& Insertion} : orientation, stages, insertion professionnelle \item \textbf{Communication \& Vie Institutionnelle} : communication interne, vie institutionnelle \item \textbf{Vie Pratique} : logement, restauration, transport, vie quotidienne \end{itemize}

\clearpage

\subsection{Règles d'Or pour la Qualité des Données}

\begin{important}[10 Règles d'Or]
\begin{enumerate}
    \item \textbf{Formulez comme un utilisateur réel} : Écrivez les questions comme un étudiant les poserait naturellement
    
    \item \textbf{Soyez spécifique} : "Quels sont les frais d'inscription en L1 ?" plutôt que "Frais L1"
    
    \item \textbf{Incluez des variantes} : Minimum 2-3 formulations différentes par réponse
    
    \item \textbf{Utilisez un vocabulaire riche} : Synonymes, expressions courantes, jargon étudiant
    
    \item \textbf{Respectez la structure CSV} : Ne jamais omettre une colonne, même si "À compléter"
    
    \item \textbf{Vérifiez la source} : Toujours indiquer l'origine de l'information
    
    \item \textbf{Réponses concises mais complètes} : Direct et informatif, sans superflu
    
    \item \textbf{Testez vos questions} : Posez-vous "comment un étudiant demanderait-il cela ?"
    
    \item \textbf{Évitez les codes} : "Comment s'inscrire au club ?" plutôt que "FAQ\_CLUB\_001"
    
    \item \textbf{Mettez à jour régulièrement} : Les informations évoluent (présidents de clubs, dates, etc.)
\end{enumerate}
\end{important}

\subsection{Checklist de Validation d'une Entrée CSV}

Avant d'ajouter une nouvelle ligne à votre fichier CSV, vérifiez :

\begin{enumerate}
    \item ☐ La question est-elle formulée en langage naturel ?
    \item ☐ Contient-elle au moins 5 mots significatifs ?
    \item ☐ Ai-je créé 2-3 variantes de cette question ?
    \item ☐ La réponse est-elle précise et vérifiée ?
    \item ☐ La catégorie et le sous-thème sont-ils corrects ?
    \item ☐ La source est-elle indiquée ?
    \item ☐ Aucune colonne n'est vide (sauf si "À compléter")
    \item ☐ Le format CSV est-il respecté (virgules, guillemets si nécessaire)
    \item ☐ Ai-je relu pour détecter les fautes d'orthographe ?
    \item ☐ Un utilisateur réel poserait-il la question ainsi ?
\end{enumerate}

\clearpage

% Section 5 : Cas d'Utilisation
\section{Cas d'Utilisation du Système}
\label{sec:cas_utilisation}

\subsection{CU-01 : Question Simple avec Bonne Correspondance}

\textbf{Acteur :} Étudiant actuel ou candidat

\textbf{Objectif :} Obtenir une réponse rapide à une question fréquente

\textbf{Scénario nominal :}

\begin{enumerate}
    \item L'utilisateur accède à l'interface chatbot (web ou mobile)
    \item Le système affiche un message de bienvenue : "Bonjour ! Comment puis-je vous aider ?"
    \item L'utilisateur tape : "Comment s'inscrire au club informatique ?"
    \item L'utilisateur clique sur "Envoyer"
    \item Le système appelle l'API chatbot avec la question
    \item Le module chatbot calcule les scores TF-IDF
    \item La meilleure correspondance est trouvée (score = 92\%)
    \begin{itemize}
        \item Question FAQ : "Quels sont les documents nécessaires pour adhérer au club informatique ?"
        \item Réponse : "Pour adhérer au club informatique, il suffit d'être étudiant inscrit à SUP'PTIC."
    \end{itemize}
    \item Le système affiche la réponse en moins de 2 secondes
    \item Le système propose 2 questions similaires :
    \begin{itemize}
        \item "Comment s'inscrire à un club cette année ?" (85\%)
        \item "Peut-on s'inscrire à plusieurs clubs ?" (72\%)
    \end{itemize}
\end{enumerate}

\textbf{Postconditions :} L'utilisateur a obtenu sa réponse et peut explorer d'autres questions liées.

\textbf{Variantes :}
\begin{itemize}
    \item Si le score est entre 60-80\%, afficher un avertissement : "Voici la réponse la plus proche :"
    \item Si score < 60\%, afficher : "Je n'ai pas trouvé de réponse précise. Voici les catégories disponibles..."
\end{itemize}

\subsection{CU-02 : Question Ambiguë Nécessitant Clarification}

\textbf{Acteur :} Nouvel étudiant

\textbf{Objectif :} Obtenir de l'aide malgré une formulation imprécise

\textbf{Scénario nominal :}

\begin{enumerate}
    \item L'utilisateur tape une question courte : "Durée formation"
    \item Le système calcule les scores et trouve plusieurs correspondances moyennes :
    \begin{itemize}
        \item "Quelle est la durée de formation pour devenir Ingénieur des Télécommunications ?" (68\%)
        \item "Quelle est la durée de formation pour devenir Technicien des Télécommunications ?" (67\%)
        \item "Quelle est la durée de formation au premier cycle management ?" (65\%)
    \end{itemize}
    \item Le système affiche : "Votre question est un peu vague. Voulez-vous savoir :"
    \item Liste les 3 meilleures correspondances comme options cliquables
    \item L'utilisateur clique sur "Ingénieur des Télécommunications"
    \item Le système affiche la réponse complète : "La formation dure 2 ans."
\end{enumerate}

\textbf{Postconditions :} L'utilisateur a clarifié sa question et obtenu la bonne réponse.

\textbf{Importance pour les données CSV :} Ce cas montre pourquoi nous avons besoin de questions spécifiques dans notre CSV, pas juste "Durée formation".

\subsection{CU-03 : Question Hors Périmètre}

\textbf{Acteur :} Visiteur externe

\textbf{Objectif :} Poser une question non couverte par notre FAQ

\textbf{Scénario nominal :}

\begin{enumerate}
    \item L'utilisateur tape : "Où puis-je acheter des livres de physique quantique ?"
    \item Le système calcule les scores de similarité
    \item Toutes les correspondances ont un score < 40\%
    \item Le système affiche :
    \begin{quote}
    "Désolé, je n'ai pas trouvé de réponse pertinente à votre question. \\
    Je suis spécialisé dans les informations sur SUP'PTIC : \\
    - Programmes académiques \\
    - Inscriptions et admissions \\
    - Vie étudiante (clubs, associations) \\
    - Services administratifs \\
    \\
    Pour d'autres questions, vous pouvez contacter : \\
    Email: info@supptic.cm \\
    Téléphone: +237 XXX XXX XXX"
    \end{quote}
    \item L'utilisateur peut choisir une catégorie ou reformuler
\end{enumerate}

\textbf{Postconditions :} L'utilisateur comprend les limites du système et sait comment obtenir de l'aide humaine.

\section{Points Clés à Retenir}

\begin{important}[Résumé des Concepts Essentiels]
\begin{enumerate}
    \item \textbf{TF-IDF transforme les mots en nombres} en tenant compte de leur fréquence dans la phrase (TF) et de leur rareté dans le corpus (IDF)
    
    \item \textbf{La similarité cosinus mesure l'angle} entre deux vecteurs pour déterminer leur proximité sémantique
    
    \item \textbf{La qualité des questions CSV détermine directement} la performance du chatbot - une question mal formulée ne sera jamais trouvée
    
    \item \textbf{Les variantes de questions sont cruciales} pour couvrir les différentes façons dont les utilisateurs posent la même question
    
    \item \textbf{Le vocabulaire riche améliore la couverture} - utiliser des synonymes et expressions courantes augmente les chances de correspondance
    
    \item \textbf{L'architecture modulaire} (FAQ, Chatbot, Users) facilite la maintenance et l'évolution du système
    
    \item \textbf{Le flux de données est simple} : Question $\rightarrow$ TF-IDF $\rightarrow$ Similarité $\rightarrow$ Meilleure réponse
\end{enumerate}
\end{important}



\end{document}